%% ma: L12-B03-0003
%% lop: 12
%% chuong: 1
%% bai: 3
%% dang: 12.3.2
%% loai: TN4
%% muc: NB
%% dap_an: "B"
%% nguon: "(NBV)-ĐỀ SỐ 2-MỖI NGÀY 1 ĐỀ THI 2025.pdf (D02, MỖI NGÀY 1 ĐỀ - Đề số 2), Phần I Câu 2"
%% kien_thuc: "Bài 3 (tiệm cận đứng của đồ thị hàm phân thức bậc nhất trên bậc nhất)"
%% trang_thai: cho-duyet
%% ly_do: "Dạng toán mới đề xuất 12.3.2 (tiệm cận đứng của hàm phân thức, tính bằng đại số), chờ giáo viên duyệt."
\begin{ex}
Tiệm cận đứng của đồ thị hàm số $y=\dfrac{x+1}{2x-2}$ là đường thẳng
\choice{$x=-1$}{\True $x=1$}{$y=\dfrac12$}{$y=-\dfrac12$}
\loigiai{Mẫu $2x-2=0\Leftrightarrow x=1$, tại đó tử bằng $2\neq0$ và $\displaystyle\lim_{x\to1^+}y=+\infty$, $\displaystyle\lim_{x\to1^-}y=-\infty$. Vậy tiệm cận đứng là $x=1$.}
\end{ex}
